\documentclass[12pt,a4paper]{article}

% Encoding (must be before other packages)
\usepackage[utf8]{inputenc}
\usepackage[T1]{fontenc}

% Layout
\usepackage[margin=2.5cm, headheight=15pt]{geometry}
\usepackage{graphicx}
\usepackage{booktabs}
\usepackage{longtable}
\usepackage{tabularx}
\usepackage{array}
\usepackage{xcolor}
\usepackage{hyperref}
\usepackage{enumitem}
\usepackage{titlesec}
\usepackage{fancyhdr}
\usepackage{pgfgantt}
\usepackage{parskip}
\usepackage[protrusion=true,expansion=false]{microtype}
\usepackage{listings}
\usepackage{caption}
\usepackage{amsmath}
\usepackage{ragged2e}

% Colors
\definecolor{iiitbblue}{RGB}{0,62,126}
\definecolor{lightgray}{RGB}{245,245,245}
\definecolor{darkgray}{RGB}{80,80,80}
\definecolor{codegreen}{rgb}{0,0.6,0}
\definecolor{codepurple}{rgb}{0.58,0,0.82}

% Hyperref — no colored text, no visible borders, no anchor dots
\hypersetup{
    colorlinks=false,
    pdfborder={0 0 0},
    pdftitle={PE Report -- Shikhar Mutta},
    pdfauthor={Shikhar Mutta}
}

% Section formatting
\titleformat{\section}{\large\bfseries\color{iiitbblue}}{\thesection.}{0.6em}{}[\titlerule]
\titleformat{\subsection}{\normalsize\bfseries\color{darkgray}}{\thesubsection}{0.6em}{}
\titleformat{\subsubsection}{\normalsize\itshape}{\thesubsubsection}{0.6em}{}

% Header / Footer
\pagestyle{fancy}
\fancyhf{}
\lhead{\small\color{darkgray} PE Report -- CrewLink 2.0}
\rhead{\small\color{darkgray} Shikhar Mutta | MT2025114}
\cfoot{\thepage}
\renewcommand{\headrulewidth}{0.4pt}

% Listings style
\lstset{
  backgroundcolor=\color{lightgray},
  basicstyle=\ttfamily\small,
  keywordstyle=\color{iiitbblue}\bfseries,
  commentstyle=\color{codegreen},
  stringstyle=\color{codepurple},
  breaklines=true,
  frame=single,
  rulecolor=\color{darkgray},
  captionpos=b
}

% Column types
\newcolumntype{L}[1]{>{\RaggedRight\arraybackslash}p{#1}}
\newcolumntype{C}[1]{>{\Centering\arraybackslash}p{#1}}

%==============================================================================
\begin{document}

%------------------------------------------------------------------------------
% Title Page
%------------------------------------------------------------------------------
\begin{titlepage}
\begin{center}

\vspace*{1.5cm}

% Institution header
{\large\bfseries International Institute of Information Technology Bangalore}\\[0.5em]
{\normalsize Project Elective (PE)}\\[0.2em]
{\small M.Tech 2025\,--\,2027 \textbar\ Semester II}\\[1.5em]

{\color{darkgray}\rule{\linewidth}{0.6pt}}\\[2.5em]

% Main title
{\LARGE\bfseries Project Elective (PE) Report}\\[0.7em]
{\large\itshape CrewLink 2.0}\\[2em]

% Decorative separator with dot
{\color{darkgray}\rule{0.44\linewidth}{0.4pt}}\enspace
{\color{iiitbblue}$\bullet$}\enspace
{\color{darkgray}\rule{0.44\linewidth}{0.4pt}}\\[3em]

% Student info
{\normalsize\bfseries Shikhar Mutta (MT2025114)}\\[0.6em]
{\normalsize shikhar.mutta@iiitb.ac.in}\\[0.6em]
{\normalsize Guide:\ Prof.\ V N Muralidhara}\\[0.6em]
{\normalsize Schedule Optimization Lab, R-303}\\[0.6em]
{\normalsize 30 May 2026}\\

\vfill

% IIITB Logo
\includegraphics[width=4cm]{IIITB_logo.png}

\vspace{3cm}

\end{center}
\end{titlepage}
\pagenumbering{roman}
\tableofcontents
\newpage
\pagenumbering{arabic}

%==============================================================================
\section{Abstract}

This report documents the contributions made by Shikhar Mutta (MT2025114)
during the Project Elective (PE) -- Semester II at IIIT Bangalore, under the
supervision of Prof.\ V N Muralidhara, in the Schedule Optimization Lab
(R-303). The work was carried out on the \textbf{CrewLink 2.0} system -- a
web-based crew rostering and scheduling optimization platform developed for
Indian Railways crew management.

Over the course of Semester II (\textbf{1 January 2026 -- 29 May 2026}),
contributions spanned three feature areas: (1)~Train Timetable UI improvements
(Crew-25/Crew-89 issues), (2)~the Copy CrewLink feature (full-stack), and
(3)~CSV Import enhancement with multi-stage conflict detection and resolution,
plus a soft-delete bug fix. In total, \textbf{11 commits} were made to the
branch \texttt{shikhar-2025-26}, contributing approximately \textbf{900+ net
lines of code} across frontend (AngularJS) and backend (Java/Spring Boot)
components.

%==============================================================================
\section{Introduction}

\subsection{Background}

Railway crew scheduling is a complex combinatorial optimization problem.
The Indian Railways network operates thousands of trains daily, requiring
precise assignment of locomotive pilots, assistant loco pilots, and crew
supervisors to driving duties. Manual scheduling is error-prone and
time-intensive.

The project aims to automate and optimize this process through a software platform
that handles:
\begin{itemize}[leftmargin=1.5em]
    \item CrewLink definitions (crew availability windows at stations)
    \item Round-trip planning for crew cycles
    \item Driving duty and driving section management
    \item Train timetable ingestion and visualization
    \item Plan generation using optimization algorithms
\end{itemize}

\subsection{Project: CrewLink 2.0}

CrewLink 2.0 is a full-stack web application built on:
\begin{itemize}[leftmargin=1.5em]
    \item \textbf{Backend:} Java 8, Spring Boot 1.3, Spring Data JPA, Spring Security
    \item \textbf{Frontend:} AngularJS (1.x), UI-Bootstrap, ocLazyLoad
    \item \textbf{Database:} MySQL 5.x (\texttt{crewlink\_db})
    \item \textbf{Server Port:} 6060
\end{itemize}

The system exposes a REST API consumed by the AngularJS Single Page Application
(SPA). Core entities include \texttt{CrewLinkVM}, \texttt{RoundTrip},
\texttt{DrivingDuty}, \texttt{DrivingSection}, \texttt{UserPlan}, and
\texttt{Train}.

\subsection{Objectives of PE}

The primary objectives assigned during PE Semester II were:
\begin{enumerate}[leftmargin=1.5em]
    \item Improve the Train Timetable UI for better usability and readability.
    \item Implement the \textit{Copy CrewLink} feature to allow duplication of
          crew scheduling plans.
    \item Enhance the CSV import pipeline with robust conflict detection and
          resolution.
    \item Fix bugs related to soft-delete behaviour in crew management.
\end{enumerate}

%==============================================================================
\section{System Overview}

\begin{table}[h]
\centering
\begin{tabular}{@{} L{3cm} L{5cm} L{6cm} @{}}
\toprule
\textbf{Layer} & \textbf{Technology} & \textbf{Role} \\
\midrule
Frontend        & AngularJS + UI-Bootstrap          & SPA Views, Controllers, Services \\
REST API        & Spring MVC + Spring Data REST     & \texttt{/api} endpoints \\
Service Layer   & Java Spring Services              & Business Logic, Optimization \\
Repository      & Spring Data JPA + Custom Impl     & DB Queries, Transactions \\
Database        & MySQL (\texttt{crewlink\_db})     & Persistence \\
\bottomrule
\end{tabular}
\caption{CrewLink 2.0 -- Architecture Layers}
\end{table}

Key domain concepts:
\begin{description}[leftmargin=1.5em]
    \item[\textbf{CrewLink}] A crew availability window at a station, defined
          by start/end time, duration, distance, and rest requirements.
    \item[\textbf{RoundTrip}] A complete crew duty cycle from home station, out
          to a remote station, and back, containing multiple Driving Duties.
    \item[\textbf{UserPlan}] A versioned crew scheduling plan. Multiple plans
          can exist; one is active at a time.
    \item[\textbf{DrivingDuty / DrivingSection}] Subdivisions of a round trip
          representing individual train segments driven by a crew member.
\end{description}

%==============================================================================
\section{Work Done -- Detailed Contributions}

\subsection{Phase 1: Train Timetable UI Improvements (February 9--10, 2026)}

\textbf{Issue Reference:} Crew-25 / Crew-89 \\
\textbf{Files Modified:} \texttt{trainTimeTable.js}, \texttt{trainTimeTable.html}

The Train Timetable view is the primary interface for viewing scheduled train
runs. Several usability issues were identified and resolved:

\subsubsection{P-1: Pagination Removal / Default to 100 Rows}

\textbf{Problem:} The timetable was paginated with a small default page size,
forcing users to navigate multiple pages to view a full day's timetable.

\textbf{Solution:} Modified the AngularJS controller (\texttt{trainTimeTable.js})
to set the default page size to 100, effectively displaying the entire
timetable in a single view for most divisions.

\subsubsection{P-3: Train Number with Day of Origination}

\textbf{Problem:} Train numbers did not indicate whether a train runs DAILY or
on specific days (MON, TUE, WED, THU, FRI, SAT, SUN).

\textbf{Solution:} Updated \texttt{trainTimeTable.html} to display the
originating day tag alongside the train number (e.g., \textit{12345/DAILY},
\textit{12346/MON}), improving crew planning accuracy.

\subsubsection{P-4 \& P-5: New Train Number Field + Editable Box}

\textbf{Problem:} There was no mechanism to add a new train number or edit an
existing one directly from the timetable UI.

\textbf{Solution:} Added a new train number input field and an inline edit box
to the timetable interface, allowing operators to update train numbers without
navigating to a separate form.

\subsubsection{P-6: Remove Day Column}

\textbf{Problem:} A redundant ``Day'' column cluttered the timetable grid.

\textbf{Solution:} Removed the column from \texttt{trainTimeTable.html} to
simplify the display, since day information was now embedded in the train
number label (P-3).

\subsubsection{Confirmation Modal}

A generic \textbf{Confirmation Modal} component was added (February 10, 2026)
to guard destructive or irreversible actions (e.g., deleting a crew link,
confirming plan creation). This reusable modal improved the overall UX safety
of the application.

\medskip
\begin{tabular}{@{} L{2.6cm} L{2.6cm} L{8cm} @{}}
\toprule
\textbf{Commit Hash} & \textbf{Date} & \textbf{Description} \\
\midrule
\texttt{d33dc18b} & 2026-02-09 & P-6: Remove Day column \\
\texttt{8f6ab4ed} & 2026-02-09 & P-1: Pagination set to 100 \\
\texttt{de63a0cb} & 2026-02-09 & P-3: Train number with originating day \\
\texttt{bba1f804} & 2026-02-09 & P-4: New train number field \\
\texttt{a1034b97} & 2026-02-09 & P-4,5: Train number edit box refined \\
\texttt{16820893} & 2026-02-10 & Confirmation Modal added \\
\texttt{badc0ecb} & 2026-02-10 & Minor update / polish \\
\bottomrule
\end{tabular}

%------------------------------------------------------------------------------

\subsection{Phase 2: Copy CrewLink Feature (March 20--26, 2026)}

\textbf{Files Modified/Created:}
\begin{itemize}[leftmargin=1.5em]
    \item \texttt{CrewLinkVMController.java} (+15 lines)
    \item \texttt{CrewLinkVMRepository.java} (+5 lines)
    \item \texttt{CrewLinkVMRepositoryImpl.java} (+201 lines)
    \item \texttt{CrewLinkCopyHelper.java} (new, +56 lines)
    \item \texttt{crewLink.js} (frontend, +60 lines modified)
\end{itemize}
\textbf{Total:} 315 insertions across 5 files.

\subsubsection{Motivation}

Railway divisions periodically need to replicate an existing crew schedule plan
to a new time period or division, making small adjustments rather than building
from scratch. The \textit{Copy CrewLink} feature enables this workflow.

\subsubsection{Backend Implementation}

A new helper class \textbf{CrewLinkCopyHelper} was introduced to separate the
copy business logic from the HTTP layer. The core copy routine in
\textbf{CrewLinkVMRepositoryImpl} performs a transactional deep copy of the
source \texttt{crew\_link} and all its children (\texttt{round\_trip},
\texttt{driving\_duty}, \texttt{driving\_duty\_element},
\texttt{driving\_section}) into the target \texttt{UserPlan}:

\begin{itemize}[leftmargin=1.5em]
    \item Each entity is inserted with \texttt{copy\_of} set to its source ID,
          enabling idempotent re-copy.
    \item All scheduling parameters (time windows, durations, distances,
          rest requirements, crew type) are preserved verbatim.
    \item Conflict detection (\texttt{findCopyConflicts}) identifies driving
          sections in the target plan whose (train, from-station, to-station,
          start day, start time) tuple overlaps with the source, allowing the
          operator to resolve before committing the copy.
\end{itemize}

Two REST endpoints were added to \textbf{CrewLinkVMController}: a
\texttt{GET /checkCopyConflicts} to surface conflicts, and a
\texttt{POST /copyCrewLink} that accepts the resolved conflict list and
performs the copy atomically.

\subsubsection{Frontend Implementation}

The AngularJS controller \texttt{crewLink.js} was updated to:
\begin{itemize}[leftmargin=1.5em]
    \item Add a ``Copy Plan'' button in the UI
    \item Invoke the backend copy endpoint via \texttt{\$http}
    \item Reload the crew link list after a successful copy operation
    \item Display appropriate success/error feedback to the operator
\end{itemize}

\medskip
\begin{tabular}{@{} L{2.6cm} L{2.6cm} L{8cm} @{}}
\toprule
\textbf{Commit Hash} & \textbf{Date} & \textbf{Description} \\
\midrule
\texttt{853330e8} & 2026-03-20 & Frontend scaffold for Copy CrewLink \\
\texttt{b2d3283f} & 2026-03-26 & Full backend + frontend implementation \\
\bottomrule
\end{tabular}

%------------------------------------------------------------------------------

\subsection{Phase 3: CSV Import Enhancement with Conflict Detection (May 28--29, 2026)}

\textbf{Files Modified/Created:}
\begin{itemize}[leftmargin=1.5em]
    \item \texttt{import-crewlink.js} (new standalone script, +247 lines)
    \item \texttt{CrewLinkCsvImportService.java} (major rework, +245 lines)
    \item \texttt{CrewLinkCsvImportController.java} (+7 lines)
    \item \texttt{crewLinkCsvImport.js} (frontend controller, +26 lines)
    \item \texttt{crewLinkCsvImport.html} (view, +5 lines)
    \item \texttt{crew1.csv} (sample import file, 100 rows)
\end{itemize}
\textbf{Total:} 590 insertions across 7 files.

\subsubsection{Motivation}

The existing CSV import for crew data lacked validation. When imported crew
links conflicted with existing data (duplicate station pairs, overlapping time
windows), the system would either silently overwrite records or crash. A robust
conflict detection and resolution mechanism was required.

\subsubsection{Conflict Detection Logic}

The reworked \textbf{CrewLinkCsvImportService} implements a multi-stage import
pipeline:

\begin{enumerate}[leftmargin=1.5em]
    \item \textbf{Parse:} Read and deserialize CSV rows using OpenCSV into
          intermediate DTOs.
    \item \textbf{Validate:} Check each row for mandatory fields, valid station
          codes, and time format compliance (\texttt{HH:MM}).
    \item \textbf{Conflict Check:} Query the database for existing crew links
          matching the same station pair and plan. If overlapping records are
          found, they are flagged.
    \item \textbf{Resolution:} Conflicting records are surfaced to the user
          with options to:
        \begin{itemize}
            \item \textit{Overwrite} -- replace the existing record
            \item \textit{Skip} -- ignore the conflicting import row
            \item \textit{Merge} -- apply field-level updates
        \end{itemize}
    \item \textbf{Commit:} Non-conflicting rows are persisted immediately;
          resolved conflicts are applied per user choice.
\end{enumerate}

\subsubsection{Frontend Changes}

The import view (\texttt{crewLinkCsvImport.html}) was updated to display
conflict summary counts and provide per-conflict resolution controls.

\subsubsection{Soft Delete Bug Fix in Copy CrewLink (May 29, 2026)}

A data-corruption bug was identified in the Copy CrewLink backend, originating
from the initial March implementation of \texttt{performCrewLinkCopy}.

\textbf{Root cause.} The original implementation used bulk \texttt{INSERT \ldots
SELECT} statements to copy all round trips for a source crew link at once, then
fetched the newly-created IDs and built an \textit{index-based} mapping from old
IDs to new IDs. Driving duties, duty elements, and driving sections were then
linked to new parents using this positional map.

This approach broke in two scenarios related to soft-deleted \texttt{UserPlan}
records:
\begin{itemize}[leftmargin=1.5em]
    \item \textbf{Partial deletion:} After conflict resolution deleted some but
          not all child entities, the remaining entities no longer occupied
          contiguous positions in the ID sequence. The positional mapping then
          wired child records to the wrong parents, producing an inconsistent
          copy.
    \item \textbf{Re-copy after soft delete:} When a \texttt{UserPlan} was
          soft-deleted (\texttt{deleted = 1} on the \texttt{user\_plan} row),
          its associated crew links and child records remained in the database.
          Re-copying the same source into a fresh plan would encounter stale
          \texttt{copy\_of} references and generate duplicate or mis-linked
          records because the bulk approach did not check for existing copies.
\end{itemize}

\textbf{Fix.} The bulk batch approach was replaced with an entity-by-entity
iterative copy in \texttt{CrewLinkVMRepositoryImpl.performCrewLinkCopy()}
(commit \texttt{600196b8}). For each level of the hierarchy the method now:
\begin{enumerate}[leftmargin=1.5em]
    \item Queries for an existing copy in the target plan using
          \texttt{copy\_of = \textlangle source\,id\textrangle}.
    \item If found -- \texttt{UPDATE} the existing record to re-link it to the
          new parent.
    \item If not found -- \texttt{INSERT} a new record with \texttt{copy\_of}
          set to the source ID.
\end{enumerate}
This makes the copy operation idempotent and correct under any combination of
prior conflict deletions or soft-deleted plans.

\medskip
\begin{tabular}{@{} L{2.6cm} L{2.6cm} L{8cm} @{}}
\toprule
\textbf{Commit Hash} & \textbf{Date} & \textbf{Description} \\
\midrule
\texttt{e86de557} & 2026-05-28 & CSV import with conflict detection \& resolution \\
\texttt{600196b8} & 2026-05-29 & Soft delete related issue resolved \\
\bottomrule
\end{tabular}

%==============================================================================
\section{Timeline of Work}

\subsection{Summary Table}

\begin{longtable}{@{} L{2.5cm} L{3cm} L{8cm} @{}}
\toprule
\textbf{Period} & \textbf{Feature Area} & \textbf{Tasks Completed} \\
\midrule
\endfirsthead
\toprule
\textbf{Period} & \textbf{Feature Area} & \textbf{Tasks Completed} \\
\midrule
\endhead
\midrule
\multicolumn{3}{r}{\small\itshape Continued on next page\ldots}
\endfoot
\bottomrule
\endlastfoot

Jan 2026
    & Onboarding \& Study
    & Codebase walkthrough; understanding domain entities (CrewLink, RoundTrip,
      DrivingDuty); local environment setup; review of open issues
      (Crew-25, Crew-89). \\[0.4em]

Feb 2026 (09--10)
    & Timetable UI Fixes
    & P-1: Pagination set to 100; P-3: Day-of-origination label; P-4/P-5:
      Train number edit box; P-6: Day column removed; Confirmation Modal. \\[0.4em]

Mar 2026 (20--26)
    & Copy CrewLink
    & Frontend scaffold; backend logic via \texttt{CrewLinkCopyHelper} +
      repository; REST endpoint; AngularJS frontend integration. \\[0.4em]

Apr 2026
    & Review \& Design
    & End-to-end testing of Copy CrewLink with real plans; design of
      multi-stage CSV conflict detection strategy. (No code commits.) \\[0.4em]

May 2026 (28--29)
    & CSV Import \& Bug Fix
    & Enhanced import pipeline; multi-stage conflict detection \& resolution
      engine; soft-delete bug in Copy CrewLink fixed (position-based ID mapping
      replaced with \texttt{copy\_of}-tracked iterative copy). \\[0.4em]
\end{longtable}

\subsection{Gantt Chart}

\begin{figure}[h]
\centering
\begin{ganttchart}[
    hgrid,
    vgrid,
    x unit=2.2cm,
    y unit chart=0.70cm,
    bar/.style={fill=iiitbblue!70, rounded corners=2pt},
    bar incomplete/.style={fill=iiitbblue!30},
    title/.style={fill=iiitbblue!15, draw=iiitbblue!60},
    title label font=\small\bfseries,
    bar label font=\small,
    bar label node/.append style={align=right},
    group label font=\small\bfseries,
    canvas/.style={draw=darkgray!40}
]{1}{5}
    \gantttitle{Jan '26}{1}
    \gantttitle{Feb '26}{1}
    \gantttitle{Mar '26}{1}
    \gantttitle{Apr '26}{1}
    \gantttitle{May '26}{1} \\

    \ganttbar{Onboarding \& Study}{1}{1} \\
    \ganttbar{Timetable UI Fixes (Crew-25/89)}{2}{2} \\
    \ganttbar{Confirmation Modal}{2}{2} \\
    \ganttbar{Copy CrewLink -- Frontend}{3}{3} \\
    \ganttbar{Copy CrewLink -- Backend}{3}{3} \\
    \ganttbar{Review \& CSV Import Design}{4}{4} \\
    \ganttbar{CSV Import Enhancement}{5}{5} \\
    \ganttbar{Soft Delete Bug Fix}{5}{5}
\end{ganttchart}
\caption{Work Timeline -- Shikhar Mutta (Jan 2026 -- May 2026)}
\end{figure}

%==============================================================================
\section{Technical Skills Demonstrated}

\begin{center}
\begin{tabular}{@{} L{3.2cm} L{5.5cm} L{4.8cm} @{}}
\toprule
\textbf{Skill Area} & \textbf{Technology / Concept} & \textbf{Applied In} \\
\midrule
Backend Dev      & Java, Spring Boot, Spring Data JPA    & Copy CrewLink, CSV Import \\
REST API Design  & Spring MVC Controllers                & Copy CrewLink endpoint \\
ORM / Database   & Hibernate, MySQL, Transactions        & Repository implementation \\
Frontend Dev     & AngularJS, HTML5, Bootstrap           & All UI features \\
Data Import      & CSV parsing (OpenCSV), DTOs           & CSV Import pipeline \\
Conflict Logic   & Multi-stage validation pipelines      & CSV conflict detection \\
Bug Fixing       & Soft delete, query filtering          & Soft delete fix \\
Version Control  & Git (feature branching, commits)      & All contributions \\
\bottomrule
\end{tabular}
\end{center}

%==============================================================================
\section{Commit Log}

Complete list of commits by Shikhar Mutta on branch \texttt{shikhar-2025-26}:

\begin{center}
\begin{tabular}{@{} L{2.4cm} L{2.4cm} L{8.5cm} @{}}
\toprule
\textbf{Hash} & \textbf{Date} & \textbf{Commit Message} \\
\midrule
\texttt{d33dc18b} & 2026-02-09 & Crew-25/89 P-6: Day column deleted \\
\texttt{8f6ab4ed} & 2026-02-09 & Crew-25/89 P-1: Pagination set to 100 \\
\texttt{de63a0cb} & 2026-02-09 & Crew-25/89 P-3: Train number with run day \\
\texttt{bba1f804} & 2026-02-09 & Crew-25/89 P-4: New train number field \\
\texttt{a1034b97} & 2026-02-09 & Crew-25/89 P-4,5: Train number edit box \\
\texttt{16820893} & 2026-02-10 & Confirmation Modal added \\
\texttt{badc0ecb} & 2026-02-10 & Update (minor polish) \\
\texttt{853330e8} & 2026-03-20 & Copy CrewLink frontend basics done \\
\texttt{b2d3283f} & 2026-03-26 & Copy CrewLink functionality implemented \\
\texttt{e86de557} & 2026-05-28 & CSV import with conflict detection \\
\texttt{600196b8} & 2026-05-29 & Soft delete related issue resolved \\
\bottomrule
\end{tabular}
\end{center}

%==============================================================================
\section{Conclusion}

Over PE Semester II, meaningful contributions were made to the CrewLink 2.0
system across three distinct phases. The Train Timetable UI improvements
directly impacted operator usability by making the timetable readable in a
single view with better train identification. The Copy CrewLink feature
introduced a productivity enhancement for scheduling planners, reducing manual
rework, and a subsequent soft-delete bug fix corrected a data-corruption issue
in the copy routine where positional ID mapping produced incorrect parent-child
associations after partial conflict deletions. The CSV import enhancement with
conflict resolution is the most technically substantial contribution,
introducing a data integrity safeguard that prevents incorrect crew data from
entering the system.

The work provided practical exposure to full-stack enterprise Java development,
optimization system internals, and the domain of railway crew scheduling --
directly relevant to the project's goals under Prof.\ V N Muralidhara
at the Schedule Optimization Lab (R-303), IIIT Bangalore.

%==============================================================================
\section*{Declaration}

I hereby declare that the work presented in this report is my own, carried out
under the guidance of Prof.\ V N Muralidhara, as part of the Project Elective
course (Semester II) at IIIT Bangalore during the academic year 2025\,--\,26.

\vspace{4em}
\begin{flushright}
\textbf{Prof.\ V N Muralidhara}\\[0.3em]
Project Elective Guide\\[2em]
\textit{Date:}~\underline{\hspace{4cm}}
\end{flushright}

\end{document}
